\chapter{과제 2: 데이터와 분포}
\label{ch:hw2}

\begin{aside}
결과 자체는 \texttt{hw2/report/}에 있다. (\texttt{hw2/report\_submit/}은 채점 제출용
빈 템플릿이라 답이 들어 있지 않으므로, 이 장은 \texttt{hw2/report/}를 유일한 출처로
쓴다. 두 파일의 차이는 \Cref{sec:hw2-learned}에 적는다.) 이 장에는 \textbf{판단의 과정}을
남긴다. 무엇을 망설였고, 왜 그렇게 정했는지를 적는다. 모든 수치는
\texttt{hw2/code/p1.py}, \texttt{hw2/code/p2.py}를 다시 실행해 재확인했다.
\end{aside}

\section{데이터 세 가지}
\label{sec:hw2-data}

세 데이터셋의 크기와 성질이 이미 어떤 모델이 맞을지 힌트를 준다.

\begin{itemize}
  \item \textbf{공항 노선 수(Airport Routes)}: 공항 3409곳, 노선 수는 정수, 최소 1에서
    최대 915. 평균 19.845, 중앙값 4로 평균이 중앙값의 약 5배다. 평균과 중앙값이 크게
    벌어지는 것은 오른쪽 꼬리가 길다는 신호다(\Cref{sec:skew-sign}).
  \item \textbf{영화 평점(Movie Votes)}: 영화 4392편, 평점은 0.1 단위로 반올림된 실수,
    범위 1.9에서 8.5. 평균 6.227, 중앙값 6.3으로 둘이 거의 같다. 대칭에 가깝다는 신호다.
  \item \textbf{배우 나이(Actor Age)}: 원본 460{,}718개 생년월일 중 정제 후 458{,}015명
    (\Cref{sec:hw2-age}에서 정제 과정을 설명한다). 연 단위 정수, 15에서 99.
\end{itemize}

\section{모수 추정: 네 가지 분포, 하나의 원리}
\label{sec:hw2-estimators}

\begin{inbrief}
  \item[Why] 데이터가 어떤 분포를 따르는지 주장하려면 먼저 후보 분포마다 모수(Parameter)를
    데이터에서 뽑아내야 비교할 대상이 생긴다.
  \item[When] 후보 분포의 형태(멱함수, 지수, 균등, 정규)를 고정하고 그 안의 숫자만
    추정할 때.
  \item[Where] 뒤에 나오는 시각적 비교(\Cref{sec:hw2-compare-how})와 적합도 판정
    전 단계. 추정이 틀리면 비교 자체가 무의미하다.
  \item[How] 균등·지수·멱법칙 세 분포는 \term{최대우도추정}(Maximum Likelihood
    Estimation, MLE)을 그대로 쓴다. 정규분포만 $\sigma$를 $n-1$로 나눈 표본표준편차로
    바꿔 쓴다(자유도 보정). 이 $\sigma$를 뺀 나머지 모수는 모두 데이터가 관측될
    확률(우도, Likelihood)을 가장 크게 만드는 값이다.
\end{inbrief}

문제는 모수 하나(예: 정규분포의 $\mu$)가 아니라 \textbf{모형 자체가 맞는지}를 물었다.
그런데 모형이 맞는지 보려면 먼저 각 모형이 이 데이터에 가장 잘 맞을 때 모수가 얼마인지
정해야 공정한 비교가 된다. 임의의 모수로 비교하면 모형이 아니라 모수를 잘못 고른
탓에 진 것처럼 보일 수 있다. 그래서 네 모형 모두 모수를 먼저 고정한 뒤에만
\Cref{sec:hw2-compare-how}에서 비교한다.

\paragraph{균등분포(Uniform Distribution)와 정규분포(Normal Distribution).}
가장 단순한 둘부터 본다. 균등분포 $[a, b]$의 MLE는 관측값의 최솟값과 최댓값이다.
정규분포(\Cref{sec:normal})의 $\mu$는 MLE(표본평균, Sample Mean) $\hat\mu = \bar x$를
쓴다. $\sigma$는 MLE(즉 $n$으로 나눈 값) 대신 $n-1$로 나눈 \term{표본표준편차}(Sample
Standard Deviation)를 썼다. $n$으로 나누면 분산을 과소추정하기 때문에(자유도 보정,
Bessel's Correction) \textbf{분산}의 불편추정량(Unbiased Estimator) $s^2$를 택한 것이다.
표준편차 $s$ 자체는 여전히 약간 편향되어 있다. $n=3409, 4392$처럼 $n$이 큰 경우
두 값의 차이는 소수 셋째 자리 이하다.

\paragraph{지수분포(Exponential Distribution).}
지수분포(\Cref{sec:exponential})는 평균 대기 시간이 $1/\lambda$인 분포다. 로그우도
(Log-Likelihood)를 $\lambda$에 대해 미분해 0으로 놓으면 $\hat\lambda$가 나온다.
\begin{align*}
  \ell(\lambda) &= n\ln\lambda - \lambda\sum_i x_i   && \text{로그우도}\\
  \frac{d\ell}{d\lambda} &= \frac{n}{\lambda} - \sum_i x_i  && \text{$\lambda$로 미분}\\
  0 &= \frac{n}{\hat\lambda} - \sum_i x_i  && \text{기울기가 0인 점}\\
  \hat\lambda &= \frac{n}{\sum_i x_i} = \frac{1}{\bar x}   && \text{정리}
\end{align*}
즉 표본평균의 역수 하나면 끝난다.

\begin{example}[공항 데이터에 지수분포 MLE 적용]
\label{ex:hw2-exp-mle}
공항 노선 수의 평균은 $\bar x = 19.845$다.
\begin{align*}
  \hat\lambda &= 1/\bar x                 && \text{MLE 공식}\\
              &= 1/19.845                 && \text{평균 대입}\\
              &= 0.0504                   && \text{반올림}
\end{align*}
$\lambda$가 작을수록 대기 시간(여기서는 노선 수)의 기댓값이 크다.
\end{example}

\paragraph{멱법칙(Power Law).}
지수분포는 노선 500개짜리 공항을 사실상 불가능하다고 본다($\lambda=0.05$면 확률이
약 $e^{-25}$). 그런데 실제로 500개가 넘는 공항이 5곳 있다(\Cref{sec:hw2-airport}).
꼬리가 지수보다 느리게 줄어드는 모형이 필요하다.

\term{멱법칙}(Power Law) 분포는 $x \geq x_{\min}$에서 $f(x) \propto x^{-\alpha}$
형태의 확률밀도함수를 갖는다. 지수분포처럼 큰 값일수록 드물지만, 꼬리가 지수적으로
줄지 않고 \textbf{거듭제곱으로만} 줄어들어 훨씬 천천히 감소한다. 네트워크의 연결 수,
도시 인구, 부의 분포처럼 소수가 대부분을 차지하는 데이터에서 흔히 등장한다.

\begin{pitfall}[책에 없는 분포]
멱법칙은 \Cref{ch:distributions}에 아직 절이 없다. 여기서 처음 정의한다.
\end{pitfall}

연속형 멱법칙의 MLE는 Clauset 등(2009)의 공식을 쓴다.
\[
  \hat\alpha = 1 + \frac{n}{\sum_{i} \ln(x_i / x_{\min})}
\]
$x_{\min}$은 멱법칙이 시작되는 하한이다. 이 과제에서는 $x_{\min} = \min_i x_i$로
정했다(데이터의 최솟값 자체를 하한으로 삼는 가장 단순한 선택). 이 선택이 공항
데이터에서 $x_{\min}=1$, 영화 데이터에서 $x_{\min}=1.9$를 준다.

\begin{example}[공항 데이터에 멱법칙 MLE 적용, 축약]
\label{ex:hw2-powerlaw-mle}
$n = 3409$, $x_{\min} = 1$일 때
\begin{align*}
  \sum_i \ln(x_i/x_{\min}) &= \sum_i \ln x_i    && x_{\min}=1\text{이므로 로그가 그대로 } \ln x_i \\
                            &= 5569.4            && \text{코드로 합산한 값}\\
  \hat\alpha &= 1 + n/\sum_i \ln(x_i/x_{\min})   && \text{공식}\\
             &= 1 + 3409/5569.4                  && \text{값 대입}\\
             &= 1 + 0.612                        && \text{나눗셈}\\
             &= 1.612                            && \text{반올림}
\end{align*}
\end{example}

\Cref{tab:hw2-params}는 네 분포 모두를 두 데이터셋에 한 번에 정리한다.

\begin{table}[htbp]
\centering
\begin{tabular}{lrr}
\toprule
모수(Parameter) & 공항 노선(Airport) & 영화 평점(Movie) \\
\midrule
$\alpha$ (멱법칙 Power Law, $x_{\min}$ 포함) & $1.612$ ($x_{\min}=1$) & $1.851$ ($x_{\min}=1.9$) \\
$\lambda$ (지수 Exponential)      & $0.0504$               & $0.1606$ \\
$[a, b]$ (균등 Uniform)           & $[1, 915]$             & $[1.9, 8.5]$ \\
$\mu, \sigma$ (정규 Normal, $n{-}1$) & $19.845,\ 53.506$   & $6.227,\ 0.893$ \\
\bottomrule
\end{tabular}
\caption{두 데이터셋에서 추정한 네 분포의 모수. $n = 3409$(공항), $n = 4392$(영화).}
\label{tab:hw2-params}
\end{table}

\section{네 모형을 비교하는 방법}
\label{sec:hw2-compare-how}

\begin{inbrief}
  \item[Why] 모수를 정했다고 모형이 맞는지 알 수 있는 것은 아니다. 데이터와 모형을
    같은 그림 위에 겹쳐야 어디가 맞고 어디가 틀리는지 눈으로 확인할 수 있다.
  \item[When] 후보 분포가 여럿이고 하나를 골라야 할 때.
  \item[Where] PDF, CDF, QQ 도표(\Cref{sec:cdf}) 중 과제는 하나만 요구했다. 세
    시각화가 서로 다른 정보를 준다.
  \item[How] 분포마다 데이터와 같은 개수 $n$의 난수를 뽑아, 데이터의 CCDF/PDF와
    모형 표본의 CCDF/PDF를 겹쳐 그리고, QQ 도표와 콜모고로프--스미르노프
    거리(Kolmogorov--Smirnov Distance) $D$로 요약한다.
\end{inbrief}

\paragraph{공항: CCDF를 로그-로그 축에.}
공항 노선 수는 정수이고 꼬리가 두 자릿수 넘게 퍼져 있다. 보통 히스토그램처럼 구간을
나누면(log bin) 큰 값 쪽에 관측치가 하나뿐인 구간이 많아 그래프가 듬성듬성해진다.
그래서 값별 확률 $P(X=k)$를 점으로 찍고, \term{여누적분포함수}(Complementary CDF,
CCDF) $P(X \geq x)$를 로그-로그(log-log) 축에 그렸다. 멱법칙은 로그-로그 축에서
\textbf{직선}이 된다는 성질이 있어(양변에 로그를 취하면 $\ln P(X\geq x) \approx
-(\alpha-1)\ln x + c$, 직선의 방정식) 이 축이 멱법칙 여부를 눈으로 바로 보여준다.

\paragraph{영화: 0.1 격자에 맞춘 히스토그램.}
영화 평점은 0.1 단위로 반올림된 값이다. 구간 경계가 0.1 격자와 어긋나면 히스토그램이
들쭉날쭉해진다(어떤 구간은 값 두 개를 담고, 바로 옆 구간은 하나도 못 담는 식). 그래서
구간 너비를 0.1 격자에 맞춰 PDF(히스토그램)를 그렸고, 꼬리가 좁아 로그 축 없이도
전체가 한 그림에 들어온다.

\paragraph{역변환 표본추출(Inverse Transform Sampling).}
네 분포 각각에서 데이터와 같은 개수 $n$의 난수를 뽑아야 한다. 지수·균등·정규는
\texttt{numpy}에 표본추출 함수가 있지만, 멱법칙은 없어서 직접 만들었다. 멱법칙의
CDF는
\[
  F(x) = 1 - (x/x_{\min})^{-(\alpha-1)}
\]
$u \sim \mathrm{Uniform}(0,1)$에 대해 $u = F(x)$를 $x$에 대해 풀면
\begin{align*}
  u &= 1 - (x/x_{\min})^{-(\alpha-1)}          && \text{CDF 정의}\\
  (x/x_{\min})^{-(\alpha-1)} &= 1-u            && \text{항 이항}\\
  x/x_{\min} &= (1-u)^{-1/(\alpha-1)}          && \text{양변을 } -(\alpha-1)\text{ 제곱근}\\
  x &= x_{\min}(1-u)^{-1/(\alpha-1)}           && \text{정리}
\end{align*}
$u$를 균등분포에서 뽑아 이 식에 넣으면 멱법칙을 따르는 $x$가 나온다. 난수 시드는
6330으로 고정해 그림을 재현할 수 있게 했다.

\paragraph{QQ 도표와 $D$.}
CCDF·PDF를 겹쳐 그리면 어디가 맞고 틀리는지는 보이지만, ``얼마나'' 다른지 숫자로
비교하기는 어렵다. 분포 전체를 한 번에 대조하면서도 숫자 하나로 요약할 도구가
필요하다.

\term{분위수 대 분위수 도표}(Quantile-Quantile Plot, QQ Plot)는 $x$축에 데이터
분위수, $y$축에 모형 표본 분위수를 찍는다. 점이 기준선 $y=x$ 위에 있으면 두 분포가
같은 모양이라는 뜻이고, 선에서 벗어난 방향이 어디서 어떻게 다른지 말해준다. 로그
축을 쓸 때는 표본에 음수나 0이 섞일 수 있어(정규 표본만 해당, 지수 표본은 항상
양수다) 그 값을 먼저 제외했다.
비교를 숫자 하나로도 요약했다. \term{콜모고로프--스미르노프 거리}(KS Distance)
$D$는 데이터의 경험적 CDF와 모형 \textbf{표본}의 경험적 CDF 사이 세로 격차의
최댓값이다(\Cref{sec:cdf}). 모형 자체의 이론적 CDF가 아니라 모형에서 뽑은 표본을
쓰므로, $D$는 그 표본을 뽑은 난수 시드에 약간 좌우된다. $D=0$이면 두 CDF가 동일하다.
이 $D$를 검정 통계량으로 쓰는 두 표본 KS 검정(Two-sample KS Test)은 \Cref{sec:cd-ks-two}에서 다룬다.

\section{공항 노선 수: 멱법칙, 상한이 있는}
\label{sec:hw2-airport}

\begin{figure}[htbp]
\centering
\includegraphics[width=0.95\linewidth]{p1a_0_data.png}
\caption{공항 노선 수의 경험적 분포(Empirical Distribution). 왼쪽은 값별 확률 $P(X=k)$, 오른쪽은 CCDF
$P(X\geq x)$, 둘 다 로그-로그 축. 38\%의 공항이 노선 1--2개뿐이고, 가장 큰
공항은 915개다.}
\label{fig:hw2-airport-data}
\end{figure}

\Cref{fig:hw2-airport-data}에서 왼쪽 PDF가 이미 직선에 가깝다. 이는 멱법칙의
특징이다. 38.0\%의 공항이 노선 1--2개뿐이고, 99번째 백분위수는 285개인데 최댓값은
915개다. 왜도(Skewness)는 6.0으로(\Cref{sec:skewness}) 강하게 오른쪽으로 치우쳐
있다.

\begin{figure}[htbp]
\centering
\includegraphics[width=0.95\linewidth]{p1a_1_powerlaw.png}
\caption{공항 노선 수 대 멱법칙 모형($\alpha=1.612$, $x_{\min}=1$). CCDF는 1에서
약 60개 노선까지 데이터를 따라가고, QQ 점도 그 구간에서 $y=x$ 위에 있다. 100개를
넘으면 모형은 계속 뻗어 $10^4$--$10^5$개짜리 공항을 만들어내지만, 데이터는 915에서
멈춘다. $D=0.209$.}
\label{fig:hw2-airport-powerlaw}
\end{figure}

\Cref{tab:hw2-airport-verdict}가 네 모형을 요약한다.

\begin{table}[htbp]
\centering
\begin{tabular}{lrrl}
\toprule
분포(Distribution) & 모수(Parameter) & KS 거리 $D$ & 판정(Verdict) \\
\midrule
멱법칙(Power Law)     & $\alpha=1.612,\ x_{\min}=1$ & \textbf{0.209} & 최선, 꼬리는 과대추정 \\
지수(Exponential)     & $\lambda=0.0504$            & 0.379          & 꼬리가 너무 빨리 줄어듦 \\
정규(Normal)           & $\mu=19.8,\ \sigma=53.5$   & 0.358          & 음수 36\% 허용, 구조적 실패 \\
균등(Uniform)          & $[1, 915]$                  & 0.857          & 최악, 형태 자체가 다름 \\
\bottomrule
\end{tabular}
\caption{공항 노선 수: 네 모형의 콜모고로프--스미르노프 거리(KS Distance) $D$와 판정.}
\label{tab:hw2-airport-verdict}
\end{table}

\textbf{공항 데이터는 멱법칙을 따르되, 꼬리에 상한이 있다.} 로그-로그 축에서 CCDF가
직선인 것은 멱법칙뿐이고, 데이터의 CCDF도 1에서 약 100개 노선까지 두 자릿수 넘게
직선에 가깝다. 나머지 세 모형은 모수를 조금 바꿔도 구조적으로 못 맞는다.

\begin{itemize}
  \item \textbf{균등}은 노선 1개짜리와 915개짜리를 똑같이 흔하다고 가정하지만,
    대부분의 공항은 작다.
  \item \textbf{정규}는 대칭이고 음수를 허용한다. $\sigma$가 $\mu$의 거의 3배라
    모형 질량의 36\%가 0 미만인데, 노선 수는 음수일 수 없다. 데이터는 하한 1에서
    시작해 오른쪽으로 치우쳐 있다(왜도 6.0, \Cref{sec:skew-sign}).
  \item \textbf{지수}는 모양은 맞다(작은 값이 많고 큰 값이 드물다). 그러나 꼬리가
    너무 빨리 줄어든다. $\lambda=0.05$면 노선 500개짜리 공항의 확률은 약
    $e^{-25}$인데, 실제로 500개를 넘는 공항이 5곳 있다.
\end{itemize}

멱법칙도 완벽하지는 않다. $\hat\alpha=1.61 < 2$이므로 이론적으로 평균이 발산하는
모형이고, 표본에는 $10^5$개짜리 공항이 등장한다(\Cref{fig:hw2-airport-powerlaw}).
실제로는 공항이 3409곳뿐이고 활주로·게이트·슬롯이 허브의 성장을 막는다. 그래서
데이터의 CCDF는 100개를 넘는 지점부터 아래로 꺾인다. 가장 적절한 설명은
\textbf{상한이 있는 멱법칙}이며, 이는 \term{선호적 연결}(Preferential Attachment)로
성장하는 네트워크의 전형적인 모양이다. 이미 노선이 많은 공항에 새 노선이 더 많이
붙는다는 뜻이다.

\section{영화 평점: 정규분포에 가깝다}
\label{sec:hw2-movie}

\begin{figure}[htbp]
\centering
\includegraphics[width=0.95\linewidth]{p1b_0_data.png}
\caption{영화 평점의 경험적 분포(Empirical Distribution). 왼쪽은 0.1 격자에 맞춘 히스토그램(PDF), 오른쪽은
CDF. 중앙값 6.3, 평균 6.23으로 6.3 근처에서 한 개의 봉우리를 이루며 거의 대칭이다.
왼쪽 꼬리가 약간 더 길다(왜도 $-0.48$).}
\label{fig:hw2-movie-data}
\end{figure}

\begin{figure}[htbp]
\centering
\includegraphics[width=0.95\linewidth]{p1b_4_normal.png}
\caption{영화 평점 대 정규분포 모형($\mu=6.227$, $\sigma=0.893$). 모형 표본이
데이터와 같은 종 모양을 그리고, QQ 점은 약 4.5에서 7.5까지 $y=x$ 위에 있다. 양 끝만
어긋난다. 데이터의 최저 평점이 모형 예측보다 더 낮고(왼쪽 꼬리가 더 김), 데이터는
8.5에서 끝나지만 모형은 계속된다. $D=0.055$로 네 모형 중 가장 작다.}
\label{fig:hw2-movie-normal}
\end{figure}

\begin{table}[htbp]
\centering
\begin{tabular}{lrrl}
\toprule
분포(Distribution) & 모수(Parameter) & KS 거리 $D$ & 판정(Verdict) \\
\midrule
정규(Normal)           & $\mu=6.227,\ \sigma=0.893$   & \textbf{0.055} & 최선, 양 끝만 이탈 \\
균등(Uniform)           & $[1.9, 8.5]$                  & 0.394          & 봉우리를 무시함 \\
멱법칙(Power Law)      & $\alpha=1.851,\ x_{\min}=1.9$ & 0.485          & 단조감소, 구조적 실패 \\
지수(Exponential)      & $\lambda=0.1606$              & 0.490          & 단조감소, 최악 \\
\bottomrule
\end{tabular}
\caption{영화 평점: 네 모형의 KS 거리 $D$와 판정.}
\label{tab:hw2-movie-verdict}
\end{table}

\textbf{영화 평점은 근사적으로 정규분포를 따른다.} $D=0.055$는 다음으로 좋은 모형의
약 7분의 1이다. 정규분포는 $\pm1\sigma$ 안에 68.3\%, $\pm2\sigma$ 안에 95.4\%를
예측하는데, 실제 데이터는 각각 70.1\%와 95.6\%로 거의 일치한다.

나머지 세 모형은 구조적으로 실패한다. 멱법칙과 지수분포는 둘 다
\textbf{단조감소}(monotonically decreasing) 형태라 가장 가능성이 높은 값이 가장
작은 값이다. 즉 대부분의 영화가 최악의 평점을 받고 소수만 매우 높은 평점을 받는다고
예측하는데, 실제 데이터는 6.3 근처에 봉우리 하나가 있다. 멱법칙 표본은 13.3\%가
20점을 넘어가는데, 평점은 1에서 10 사이만 가능하다. 균등분포는 봉우리 자체를
무시한다.

정규 적합도 완벽하지 않다. 첫째, 평점은 $[1, 10]$에 갇혀 있고 약간 왼쪽으로
치우쳐 있다(\Cref{sec:skew-sign}). 나쁜 영화는 평균보다 훨씬 더 낮게 갈 수 있지만
좋은 영화가 평균보다 더 높이 갈 여지는 상대적으로 작다. 둘째, 평점이 소수점
한 자리로 반올림되어 QQ 도표에 작은 계단이 생긴다.

\term{중심극한정리}(Central Limit Theorem, CLT)는 영화 한 편의 평점이 그 영화의
개별 투표 여러 개의 평균이라는 사실을 설명해 줄 수는 있지만, 그것이 설명하는 대상은
영화 \textbf{한 편} 안에서 평점이 진짜 평균 주변에 얼마나 좁게 모이는지이지, 서로
다른 영화 4392편의 평균 평점이 \textbf{영화 간}에 왜 종 모양으로 퍼지는지가 아니다.
후자는 영화마다 다른 진짜 작품성의 분포에서 온다. 개별 투표의 잡음이 평균에서
줄어드는 것과, 영화 간 평점 분포가 정규 모양을 띠는 것은 서로 다른 현상이며, 이
장에서 관찰한 것은 후자, 즉 \Cref{fig:hw2-movie-data}와 \Cref{tab:hw2-movie-verdict}가
보여주는 데이터 자체의 형태다.

\section{나이, 첫째 자리, 마지막 자리}
\label{sec:hw2-age}

\begin{pitfall}[생년월일 정제: 두 자리 연도와 결측 표시]
생년월일은 \texttt{dd-Mon-yy} 형식이 대부분이고 일부만 4자리 연도다. 모든 두 자리
연도를 \textbf{19\textit{yy}}로 읽었다. 이 데이터(IMDB-WIKI)는 2015년에 수집됐으므로
``20''--``24''는 2020--2024년일 수 없고, 연도별 개수를 세어 보면 1920년대부터
``00''까지 매끄럽게 줄어들 뿐 2000년대로 전환됐다고 볼 만한 단절이 없다. 이 확인을
거친 뒤에야 19\textit{yy} 규칙을 확정했다(짐작으로 정하지 않았다). 연도 ``0000''은
\term{결측값}(Missing Value) 표시로 보고 제외했다(129건). 나이는
\textbf{2015-01-01}(수집 연도) 기준
만 나이이고, 1에서 99 사이만 남겼다. 그 결과 \textbf{458{,}015명}이 남는다(평균
45.3, 중앙값 43, 범위 15--99). 최솟값이 15인 것은 데이터의 성질이 아니라
19\textit{yy} 규칙의 결과다. 가장 늦은 출생연도가 1999년이므로 2015년 기준
최소 나이가 15로 정해진다.
\end{pitfall}

\begin{figure}[htbp]
\centering
\includegraphics[width=0.8\linewidth]{p2_1_age.png}
\caption{2015년 기준 나이 분포($n=458{,}015$). 30대와 40대에 넓은 봉우리가
하나 있고, 99세까지 긴 오른쪽 꼬리가 이어진다. 대부분의 질량은 25에서 60 사이다.}
\label{fig:hw2-age}
\end{figure}

\begin{figure}[htbp]
\centering
\includegraphics[width=0.8\linewidth]{p2_2_first_digit.png}
\caption{나이의 첫째 자리. 점선은 균등 기준선 $1/9$, 주황 곡선은
벤포드 법칙(Benford's Law) $\log_{10}(1+1/d)$. 자릿수 1부터 9까지 비율은
1.3\%, 11.3\%, 26.2\%, 28.3\%, 16.7\%, 9.0\%, 4.4\%, 2.2\%, 0.7\%다. 데이터는
균등과도 벤포드와도 거리가 멀다. 벤포드는 1에서 꾸준히 감소하지만 데이터는 3과
4에서 정점을 찍는다.}
\label{fig:hw2-first-digit}
\end{figure}

\begin{figure}[htbp]
\centering
\includegraphics[width=0.8\linewidth]{p2_3_last_digit.png}
\caption{나이의 마지막 자리. 점선은 균등 기준선 $1/10$. 모든 자릿수가
8.95\%에서 11.01\% 사이, 즉 10\%에서 최대 1.1퍼센트포인트 안쪽에 있다.}
\label{fig:hw2-last-digit}
\end{figure}

\textbf{마지막 자리는 (근사적으로) 균등하고, 첫째 자리는 그렇지 않다.}

\paragraph{마지막 자리.} 이는 예상된 결과다. 나이의 마지막 자리는 나이를 10으로 나눈
나머지다. 20대, 30대처럼 10년 단위 구간마다 마지막 자리 0에서 9가 정확히 한 번씩
들어 있다. 나이 분포는 한 살 단위로 보면 완만하게 변하므로, 어느 10년 구간을 보든
그 안의 열 개 나이는 거의 같은 비율을 차지하고, 이를 모든 구간에 걸쳐 더하면 각
마지막 자리가 약 10\%씩 나온다. 남은 $\pm1.1$퍼센트포인트의 편차는 특정 자릿수를
선호해서가 아니라, 특정 출생연도에 배우가 몇 명 태어났는지의 연도별 잡음
(\Cref{fig:hw2-age}의 들쭉날쭉한 막대)에서 온다.

\paragraph{첫째 자리.} 이것도 사실은 예상된 결과다. 나이 10에서 99까지는 첫째
자리가 곧 \textbf{인생의 10년 단위 구간(decade)}이다. 그래서 첫째 자리의 히스토그램은
나이 분포를 10년 단위로 묶은 것과 같은 모양이고, 실제로 \Cref{fig:hw2-age}와 같은
형태(봉우리 하나 + 오른쪽 꼬리)를 보인다. 자릿수 1이 거의 비어 있는 것은 15--19세만
가능한 좁은 구간이라 그 나이대 배우가 적기 때문이고, 9가 거의 비어 있는 것은 90대에
이르는 배우가 드물기 때문이다. 첫째 자리가 균등하려면 10대부터 90대까지 모든
10년 구간에 배우 수가 같아야 하는데, 활동 중인 배우 집단은 그렇지 않다.

\paragraph{벤포드 법칙도 아니다.} \term{벤포드 법칙}(Benford's Law, 첫째 자리
$d$의 확률이 $\log_{10}(1+1/d)$이라 1이 30.1\%로 가장 흔함)은 값이 여러 자릿수
(order of magnitude)에 걸쳐 퍼져 있는 데이터, 예를 들어 인구, 가격, 또는
\Cref{sec:hw2-airport}의 공항 노선 수처럼 로그값이 고르게 퍼진 데이터에서 성립한다.
나이는 15에서 99까지 한 자릿수 범위 안에 있고 45 근처에 몰려 있어서, 첫째 자리는
벤포드의 패턴이 아니라 그 봉우리가 어디 있는지를 보여줄 뿐이다. 요약하면, 마지막
자리는 값이 10년 구간의 어디에 있는지만 남기고 나머지 정보를 버리기 때문에
균등하고, 첫째 자리는 바로 그 정보(어느 구간인가)를 그대로 담고 있어서 나이
분포의 모양을 따라간다.

\section{이 과제에서 배운 것}
\label{sec:hw2-learned}

\textbf{모형이 맞는지는 모수 하나가 아니라 꼬리의 모양으로 갈린다.} 공항은
$\sigma$(53.5)가 $\mu$(19.8)의 거의 3배라 평균과 표준편차만 봐도 이미 정규분포가
아니다(\Cref{sec:hw2-airport}). 영화는 평균과 중앙값이 거의 같고 왜도도 작아 정규분포
후보로 보였고, 실제로 그렇게 나왔다(\Cref{sec:hw2-movie}). 두 경우 모두 갈림을 만든
것은 꼬리의 모양이었다. 공항은 꼬리가 지수보다 느리게 줄어 멱법칙이 이겼고, 영화는
꼬리가 짧고 대칭이라 정규분포가 이겼다.

\textbf{KS 거리 하나로는 ``왜'' 지는지 알 수 없다.} \Cref{tab:hw2-airport-verdict}와
\Cref{tab:hw2-movie-verdict}의 $D$ 값은 순위를 매겨 줄 뿐, 그 순위의 이유는 말해주지
않는다. 공항에서는 정규($D=0.358$)가 지수($D=0.379$)보다 근소하게 \textbf{낫지만},
정규는 모형 질량의 36\%를 음수 영역에 두는 구조적 실패이고 지수는 꼬리가 너무 빨리
줄어드는 실패다. $D$가 더 작다고 더 타당한 모형은 아니다. 숫자 뒤에 그림을 반드시
같이 본 이유다.

\textbf{자릿수 분포는 원본 분포의 그림자다.} 마지막 자리가 균등한 것과 첫째 자리가
균등하지 않은 것은 서로 다른 현상이 아니라, 같은 나이 분포를 \textbf{정보를 버리는
정도가 다른 두 렌즈}로 본 것이다. 벤포드 법칙이 성립하려면 값이 여러 자릿수에
걸쳐 퍼져 있어야 한다는 조건이 있고, 나이는 그 조건을 만족하지 않는다.

마지막으로 자료 확인 하나. \texttt{hw2/report\_submit/hw2\_taeklim\_report.tex}는
채점 제출용 빈 템플릿으로, 문항마다 \texttt{STEP} 주석만 있고 그림과 답은 모두
주석 처리되어 있다. 실제 답과 논의, 그리고 그림·표는 전부
\texttt{hw2/report/hw2\_taeklim\_report.tex}에 있다. 이 장의 모든 수치와 그림은
후자를 출처로 삼았다.
